{\setlength{\parskip}{4mm}
{\huge{\textbf{Resumen}}}

El desarrollo de aplicaciones web de gestión de datos en entornos de
investigación se enfrenta a una tensión recurrente: la necesidad de entregar
software de calidad en plazos cortos frente a la inevitable evolución del
modelo de datos a lo largo del proyecto. Cada vez que el modelo cambia, el
coste de mantener sincronizados el backend, la capa de comunicación y la
interfaz de usuario crece de forma desproporcionada, ya que buena parte de ese
código se escribe a mano y se acopla a un dominio concreto.

Este Trabajo de Fin de Grado aborda dicho problema mediante el diseño,
refactorización y consolidación de un \emph{framework} \emph{full-stack}
reutilizable, compuesto por un backend en Python (\emph{uniback}, basado en
FastAPI, SQLAlchemy~2.x y PostgreSQL) y una librería de componentes Angular
(\emph{ngt-gui}). El núcleo de la contribución es una cadena de generación
automática que parte del modelo ORM, deriva un \emph{JSON Schema} enriquecido
con pistas de presentación y, a partir de él, construye formularios dinámicos
con Formly sin escribir HTML específico de cada entidad. El trabajo se
estructura en una metodología iterativa incremental y se valida sobre las
necesidades reales de los proyectos del Instituto Tecnológico de Canarias con
el Jardín Botánico Viera y Clavijo.

Los resultados muestran un contrato de comunicación uniforme entre backend y
frontend, un sistema de descubrimiento de entidades, una capa cliente
genérica tipada y un generador de formularios guiado por el esquema, todo ello
respaldado por pruebas automatizadas. El framework reduce drásticamente el
código repetitivo necesario para exponer y editar una entidad nueva, y queda
preparado para su reutilización en proyectos con modelos de datos en
evolución.

\textbf{Palabras clave:} framework, generación automática de formularios,
JSON Schema, FastAPI, Angular, Formly, modelo de datos en evolución, desarrollo
rápido de aplicaciones.

\newpage

{\huge{\textbf{Abstract}}}

Building data-management web applications in research environments faces a
recurring tension: delivering quality software on short schedules versus the
inevitable evolution of the data model throughout the project. Every time the
model changes, the cost of keeping the backend, the communication layer and
the user interface in sync grows disproportionately, because much of that code
is hand-written and coupled to a specific domain.

This Bachelor's Thesis tackles that problem through the design, refactoring and
consolidation of a reusable full-stack framework, composed of a Python backend
(\emph{uniback}, based on FastAPI, SQLAlchemy~2.x and PostgreSQL) and an
Angular component library (\emph{ngt-gui}). The core contribution is an
automatic generation pipeline that starts from the ORM model, derives an
enriched JSON Schema with presentation hints and, from it, builds dynamic
Formly forms without writing entity-specific HTML. The work follows an
iterative and incremental methodology and is validated against the real needs
of the projects carried out by the Instituto Tecnológico de Canarias with the
Jardín Botánico Viera y Clavijo.

The results include a uniform communication contract between backend and
frontend, an entity discovery system, a generic typed client layer and a
schema-driven form generator, all backed by automated tests. The framework
sharply reduces the boilerplate required to expose and edit a new entity, and
is prepared for reuse in projects with evolving data models.

\textbf{Keywords:} framework, automatic form generation, JSON Schema, FastAPI,
Angular, Formly, evolving data model, rapid application development.

}
